\chapter{Giới thiệu đề tài}
\label{ch:intro}

\section{Đặt vấn đề}

Trong những năm gần đây, cùng với sự phát triển mạnh mẽ của công nghệ thông tin và thương
mại điện tử, mô hình kinh doanh trực tuyến ngày càng trở nên phổ biến tại Việt Nam. Hầu hết
các cửa hàng và doanh nghiệp bán lẻ đã và đang chuyển đổi từ hình thức kinh doanh truyền
thống sang trực tuyến nhằm tiếp cận khách hàng nhanh hơn, tiết kiệm chi phí vận hành và nâng
cao hiệu quả kinh doanh.

Trước đây, việc quản lý sản phẩm, đơn hàng hay dữ liệu khách hàng chủ yếu được thực hiện
thủ công, gây ra nhiều hạn chế trong quá trình vận hành, đặc biệt khi quy mô kinh doanh mở
rộng. Với các mặt hàng có tính cập nhật nhanh và đa dạng mẫu mã như quần áo trẻ em, việc
quản lý thủ công càng bộc lộ nhiều bất cập --- từ khâu cập nhật sản phẩm, theo dõi tồn kho
theo từng biến thể (kích cỡ, màu sắc), đến xử lý đơn hàng và chăm sóc khách hàng.

Lý do lựa chọn ngành hàng thời trang trẻ em xuất phát từ đặc thù riêng của nhóm sản phẩm này.
Quần áo trẻ em có vòng đời ngắn do trẻ lớn nhanh, nhu cầu thay mới liên tục, phân loại theo
nhiều độ tuổi và giới tính, đồng thời đặt ra yêu cầu cao về chất liệu và độ an toàn. Người mua
chủ yếu là các bậc phụ huynh trẻ, bận rộn, đề cao sự tiện lợi cũng như thông tin sản phẩm minh
bạch và đáng tin cậy. Đây là một thị trường có sức tiêu thụ ổn định và tiềm năng lớn cho hoạt
động kinh doanh trực tuyến.

Qua quá trình khảo sát thực tế, tôi nhận thấy nhiều cửa hàng thời trang trẻ em quy mô vừa và
nhỏ vẫn đang bán hàng chủ yếu qua các kênh mạng xã hội như Facebook, Instagram hay các sàn
thương mại điện tử trung gian. Cách làm này tuy dễ bắt đầu nhưng bộc lộ nhiều hạn chế: thông
tin sản phẩm phân tán, khó quản lý tồn kho và đơn hàng tập trung, không chủ động được về
thương hiệu, dữ liệu khách hàng và các chương trình khuyến mãi. Việc xây dựng một website
thương mại điện tử riêng sẽ giúp cửa hàng giới thiệu sản phẩm chuyên nghiệp, quản lý hàng hóa
và đơn hàng hiệu quả, đồng thời mang lại trải nghiệm mua sắm trực quan, tiện lợi và hiện đại
cho người dùng.

Từ thực tế đó, tôi quyết định thực hiện đề tài \textit{``\thesisTitle''} với mục tiêu xây dựng một
nền tảng thương mại điện tử hiện đại, tối ưu trải nghiệm khách hàng, đồng thời hỗ trợ người bán
trong việc quản lý và phát triển thương hiệu trong thời đại số.

\section{Mục tiêu và phạm vi đề tài}

\subsection{Mục tiêu}

Mục tiêu hàng đầu của đề tài là xây dựng một website thương mại điện tử hoàn chỉnh, hiện đại,
thân thiện và dễ sử dụng cho thương hiệu thời trang trẻ em \brandName{}, đáp ứng đồng thời hai
nhóm nhu cầu:

\begin{itemize}
  \item \textbf{Đối với khách hàng:} cung cấp trải nghiệm mua sắm trực tuyến thuận tiện ---
  duyệt và tìm kiếm sản phẩm theo danh mục, độ tuổi, giới tính, kích cỡ và màu sắc; xem chi
  tiết sản phẩm với hình ảnh và đánh giá; quản lý giỏ hàng, danh sách yêu thích; đặt hàng và
  thanh toán trực tuyến an toàn; theo dõi trạng thái đơn hàng.

  \item \textbf{Đối với người bán (quản trị):} cung cấp một hệ quản trị tập trung cho phép quản
  lý sản phẩm và biến thể, tồn kho, đơn hàng, khách hàng, khuyến mãi (mã giảm giá, flash sale),
  nội dung trang chủ, cũng như theo dõi báo cáo doanh thu và xuất dữ liệu phục vụ ra quyết định
  kinh doanh.
\end{itemize}

Bên cạnh mục tiêu sản phẩm, đề tài còn hướng tới việc vận dụng các công nghệ web hiện đại
(Next.js, React Server Components, Prisma, PostgreSQL) để xây dựng hệ thống theo hướng có
khả năng mở rộng, dễ bảo trì, bảo đảm hiệu năng và an toàn dữ liệu.

\subsection{Phạm vi}

Hệ thống \brandName{} được xây dựng với hai nhóm tác nhân chính: \textbf{Khách hàng} (người
mua) và \textbf{Quản trị viên} (người bán/quản lý). Phạm vi chức năng của từng nhóm như sau:

\textbf{Đối với Khách hàng:}
\begin{itemize}
  \item Đăng ký, đăng nhập, đăng xuất; quản lý thông tin cá nhân và sổ địa chỉ giao hàng.
  \item Xem danh sách và chi tiết sản phẩm; tìm kiếm và lọc theo danh mục, độ tuổi, giới tính,
  kích cỡ, màu sắc, khoảng giá.
  \item Quản lý giỏ hàng, danh sách yêu thích và so sánh sản phẩm.
  \item Tạo đơn hàng, áp dụng mã giảm giá, đặt hàng và thanh toán (COD hoặc trực tuyến qua
  PayOS); theo dõi lịch sử và trạng thái đơn hàng; yêu cầu đổi/trả hàng.
  \item Đánh giá, bình luận sản phẩm sau khi mua.
\end{itemize}

\textbf{Đối với Quản trị viên:}
\begin{itemize}
  \item Quản lý danh mục, sản phẩm và biến thể sản phẩm; quản lý tồn kho.
  \item Quản lý đơn hàng (xác nhận, cập nhật trạng thái giao hàng, mã vận đơn), xử lý đổi/trả.
  \item Quản lý mã giảm giá và chương trình flash sale; quản lý banner và nội dung trang chủ.
  \item Quản lý tài khoản khách hàng; kiểm duyệt đánh giá.
  \item Xem báo cáo thống kê doanh thu, sản phẩm bán chạy và xuất báo cáo ra tệp Excel.
\end{itemize}

\textbf{Giới hạn phạm vi:} Để tập trung vào các nghiệp vụ cốt lõi của một sàn thương mại điện
tử và phù hợp với thời gian thực hiện đồ án, hệ thống có một số giới hạn về phạm vi: thanh toán
trực tuyến sử dụng một cổng duy nhất là PayOS (bên cạnh COD); việc giao vận được quản lý ở
mức trạng thái và mã vận đơn do quản trị viên cập nhật thủ công (không tích hợp sâu với các đơn
vị vận chuyển như GHN, GHTK); tìm kiếm sản phẩm sử dụng cơ chế tìm kiếm toàn văn của
PostgreSQL thay vì dịch vụ tìm kiếm ngoài.

\section{Định hướng giải pháp}

Trên cơ sở mục tiêu và phạm vi đã xác định, đề tài lựa chọn định hướng giải pháp như sau:

\begin{itemize}
  \item \textbf{Kiến trúc full-stack hợp nhất với Next.js 14:} sử dụng App Router cùng React
  Server Components và Server Actions, cho phép phần lớn nghiệp vụ được xử lý ngay trên máy chủ,
  giảm lượng mã chạy phía trình duyệt, cải thiện hiệu năng và tối ưu SEO --- yếu tố quan trọng đối
  với một website thương mại điện tử.

  \item \textbf{Lưu trữ dữ liệu với PostgreSQL trên nền Supabase:} truy cập qua Prisma ORM để
  bảo đảm an toàn kiểu dữ liệu và năng suất phát triển; tận dụng các tính năng nâng cao của
  PostgreSQL như tìm kiếm toàn văn (\texttt{unaccent}, \texttt{pg\_trgm}) và bảo mật mức dòng
  (Row-Level Security).

  \item \textbf{Thanh toán trực tuyến qua PayOS:} tích hợp cổng thanh toán PayOS với cơ chế tạo
  liên kết thanh toán, xác thực webhook và tự động hủy đơn chưa thanh toán quá hạn để hoàn trả
  tồn kho.

  \item \textbf{Triển khai trên nền tảng Vercel:} tận dụng khả năng triển khai liên tục, hàm
  serverless và tác vụ định kỳ (cron) sẵn có cho ứng dụng Next.js.
\end{itemize}

\section{Bố cục đồ án}

Nội dung đồ án được trình bày trong sáu chương với bố cục như sau:

\begin{itemize}
  \item \textbf{Chương 1 --- Giới thiệu đề tài:} trình bày bối cảnh, lý do chọn đề tài, mục tiêu,
  phạm vi, định hướng giải pháp và bố cục của đồ án.
  \item \textbf{Chương 2 --- Khảo sát và phân tích yêu cầu:} khảo sát hiện trạng, xác định các tác
  nhân và chức năng của hệ thống thông qua biểu đồ use case, đặc tả chi tiết các ca sử dụng và các
  yêu cầu phi chức năng.
  \item \textbf{Chương 3 --- Công nghệ sử dụng:} giới thiệu các công nghệ, nền tảng và công cụ
  được sử dụng để xây dựng hệ thống cùng lý do lựa chọn.
  \item \textbf{Chương 4 --- Phát triển và triển khai ứng dụng:} trình bày thiết kế kiến trúc, thiết
  kế cơ sở dữ liệu, quá trình xây dựng, kết quả giao diện, kiểm thử và triển khai hệ thống.
  \item \textbf{Chương 5 --- Các giải pháp và đóng góp nổi bật:} phân tích một số vấn đề kỹ thuật
  tiêu biểu gặp phải trong quá trình thực hiện và giải pháp tương ứng.
  \item \textbf{Chương 6 --- Kết luận và hướng phát triển:} tổng kết các kết quả đạt được và đề
  xuất các hướng phát triển trong tương lai.
\end{itemize}
